\begin{thm}\label{thm:logconvex}Assume $(b_n)$ is $q$-increasing, cf.\ Definition~{\rm \ref{thm:defqinc}}, and let $U_n$ be given by~\eqref{eq:Un}. Then
 \begin{gather*}%\label{eq:logconvexU}
U_n\le \big(q^2;q^2\big)_\infty^{-2} \sum_{k=0}^{\infty}q^{2k^2},\qquad n\ge 0.
\end{gather*}
This holds in particular if $(b_n)$ is log-convex and strictly increasing with $q=b_0/b_1<1$.
\end{thm}

\begin{proof}For $l\ge 1$ we have
\begin{gather}\label{eq:logcv-q}
\frac{b_n}{b_{n+l}}=\frac{b_n}{b_{n+1}}\frac{b_{n+1}}{b_{n+2}}\cdots\frac{b_{n+l-1}}{b_{n+l}}\le q^l,
\end{gather}
 so from equation \eqref{eq:u} we get
\begin{gather}\label{eq:uq}
u_{n+1,k}\le q^{2k} u_{n,k}+q^{2k-1}u_{n,k-1}.
\end{gather}
We claim that
\begin{gather}\label{eq:ind1*}
u_{n,k}\le q^{k^2}\prod_{j=1}^k\big(1-q^{2j}\big)^{-1}=\frac{q^{k^2}}{\big(q^2;q^2\big)_k},
\end{gather}
which is certainly true for $n=0,k\ge 0$ since $u_{0,0}=1$ and $u_{0,k}=0$ for $k\ge 1$. It is also true when $k=0$ independent of $n$ since $u_{n,0}=1$.

Assume that \eqref{eq:ind1*} holds for some $n$ and all $k$. From \eqref{eq:uq} we get for $k\ge 1$
\begin{gather*}
u_{n+1,k}\le \frac{q^{k^2+2k}}{\big(q^2;q^2\big)_k}+\frac{q^{(k-1)^2+2k-1}}{\big(q^2;q^2\big)_{k-1}}=\frac{q^{k^2}}{\big(q^2;q^2\big)_k},
\end{gather*}
which finishes the proof of \eqref{eq:ind1*}. In particular, using the notation from~\eqref{eq:q-infty}
\begin{gather*}
u_{n,k}\le \frac{q^{k^2}}{\big(q^2;q^2\big)_\infty},
\end{gather*}
hence by \eqref{eq:u1}
\begin{gather*}
U_n\le \big(q^2;q^2\big)_\infty^{-2} \sum_{k=0}^{\infty}q^{2k^2}<\infty.\tag*{\qed}
\end{gather*}\renewcommand{\qed}{}
\end{proof}

Theorem~\ref{thm:logconvex} can be slightly extended.

\begin{thm}\label{thm:evtlogconvex} Assume $(b_n)$ is eventually $q$-increasing, cf.\ Definition~{\rm \ref{thm:defqinc}}. Then the sequen\-ce~$(U_n)$ from~\eqref{eq:Un} is bounded.
\end{thm}

\begin{proof} By \eqref{eq:qinc} we see that $(b_n)$ is strictly increasing for $n\ge n_0$ and tends to infinity. For $k\ge 1$ we then have
\begin{gather*}
\max\{b_n\colon 0\le n<n_0+k\}=\max\left(\max\{b_n\colon 0\le n<n_0\}, b_{n_0+k-1}\right)
\end{gather*}
and the latter is $<b_{n_0+k}$ for $k$ sufficiently large.

By increasing $n_0$ if necessary we may therefore assume that
\begin{gather*}
b_{n_0}> \max\{b_n\colon 0\le n<n_0\}.
\end{gather*}

Defining
\begin{gather*}
q_1=\max_{0\le n<n_0} \left(\frac{b_n}{b_{n_0}}\right)^{1/(n_0-n)}, \qquad p=\max(q,q_1),
\end{gather*}
we have $0<p<1$.

Let
\begin{gather*}
\tilde{b}_n:=\begin{cases} b_{n_0}p^{n_0-n}, & 0\le n<n_0,\\
b_n, & n_0\le n.
\end{cases}
\end{gather*}
Then $\big(\tilde{b}_n\big)$ is $p$-increasing and $b_n\le\tilde{b}_n$. By Proposition~\ref{thm:increasofbn} we get $s_{2n}\le\tilde{s}_{2n}$ and for $n >n_0$ we find
\begin{gather*}
U_n=\frac{s_{2n}}{b_0^2b_1^2\cdots b_{n-1}^2}\le \frac{\tilde{b}_0^2\cdots \tilde{b}_{n_0-1}^2}{b_0^2\cdots b_{n_0-1}^2} \tilde{U}_n.
\end{gather*}
This shows the boundedness of $(U_n)$ since $(\tilde{U}_n)$ is bounded by Theorem~\ref{thm:logconvex}.
\end{proof}
